\item \textit{¿Por qué no es posible la siguiente situación?} Un aparato se diseñó de manera que un vapor inicialmente a $T = 150\text{ }^\circ\text{C}$, $P = 1.00\text{ atm}$ y $V = 0.5000\text{ m}^3$ en un pistón cilíndrico experimente un proceso donde (1) el volumen permanezca constante y la presión baje a 0.870 atm, seguido por (2) una expansión en la cual la presión es constante y el volumen aumenta a 1.00 m$^3$, y después por (3) un retorno a las condiciones iniciales. Es importante que la presión del gas nunca sea menor que 0.850 atm para que el pistón sustente una parte muy delicada y costosa del aparato. Sin tal respaldo, el delicado aparato se dañaría severamente y quedaría inútil. Cuando el diseño se convierte en un prototipo funcional, entonces opera perfectamente.
